\begin{table}[htbp]
\centering\small
\begin{tabular}{lrrrrrr}
\toprule
Design & $T$ & Audit & Joint & $4T$ & $8T$ & Available\\
\midrule
Original & 720 & 200 & 0.9340 & 0.9400 & 0.9405 & 2000/2000\\
Baseline & 1200 & 1600 & 0.9360 & 0.9470 & 0.9485 & 2000/2000\\
Longer history & 4800 & 1600 & 0.9495 & 0.9520 & 0.9520 & 2000/2000\\
Small audit & 1200 & 200 & 0.9270 & 0.9330 & 0.9340 & 2000/2000\\
Large audit & 1200 & 6400 & 0.9380 & 0.9435 & 0.9420 & 2000/2000\\
Persistent & 1200 & 1600 & 0.9325 & 0.9400 & 0.9395 & 2000/2000\\
Near rank & 1200 & 1600 & 0.2205 & 0.2210 & 0.2210 & 451/2000\\
False exogeneity & 1200 & 1600 & 0.2290 & 0.2055 & 0.2085 & 2000/2000\\
\bottomrule
\end{tabular}
\caption{All-request coverage in 2,000 fresh requests per design. Audit is the independent count per stratum. The persistent case has $\rho=.96$; near-rank probabilities are .55 and .515. Pilot arms use nested independent histories; unavailable requests count as noncoverage. These are synthetic selected-law results.}
\label{tab:v10:coverage}
\end{table}
At 720 dates the joint procedure covers .9340 and the pilot procedures
cover .9400 and .9405. The ordinary exact binomial interval for .9400 only
slightly excludes .95; the adjusted interval includes it. With 200 audit
observations per stratum at 1200 dates, coverage is .9270/.9330/.9340. All
three exact simultaneous intervals end below .95 (upper endpoints
.94372/.94900/.94988), so the small-audit shortfall persists after accounting
for all 24 marginal comparisons. The last endpoint is close to .95; its
significance should not be enlarged into a large economic distinction.
Mean SE divided by empirical SD in that design is about .939--.944.

At 4800 dates, coverage is .9495/.9520/.9520, with all ordinary and adjusted
intervals containing .95. This is evidence for this selected design, not a
uniform theorem. The high-persistence case remains .9325/.9400/.9395; the
joint procedure's exact adjusted upper endpoint is .94856. In the near-rank
case only 451 of 2,000 requests return intervals. The pilot arms cover 442
of those available requests, yet all-request coverage is only .2210. The
conditional figure .9800 does not qualify a broadly available operation.
False exogeneity gives coverage .2290/.2055/.2085 despite all requests being
numerically available.

Paired pilot-minus-joint coverage differences at the baseline are .0110 and
.0125. An exact conditional binomial test of equal coverage gives Bonferroni
adjusted two-sided $p$ values .0394 and .00408 across the 24 paired tests.
Thus there is selected-design evidence favoring the pilots over joint
regression. The conservative simultaneous intervals for the difference,
constructed by separate gain/loss bounds, still include zero; these bounds
are looser than the exact conditional test and are not its inversion.
Neither result establishes general superiority. Every conservative interval
for eightfold minus fourfold pilot coverage includes zero. Even at 4800 dates
the two pilot arms have identical covering counts; doubling pilot history
has no demonstrated coverage advantage here.

The initial normal approximation gives a zero-width paired interval when
both discordance counts are zero. That is not population certainty. Exact
binomial inversion for gain and loss probabilities, followed by a union
bound, gives a conservative difference interval $[L_g-U_l,U_g-L_l]$.
For the near-rank pilot comparison with no observed discordance, the ordinary
conservative interval is $[-.00219,.00219]$, and the simultaneous interval
is $[-.00377,.00377]$. This reporting repair leaves the procedures and draws
unchanged and retains the original approximations. Exact marginal simultaneous
intervals use error .05/24 per cell. Paired simultaneous bounds use .05/48
for each of the two binomial proportions in 24 comparisons. Each family is
controlled separately; no combined .95 guarantee across both families is
claimed. The binomial count inversion is the elementary construction in
Section~\ref{s61:binomial-inversion} applied to independent simulation requests.

\begin{table}[htbp]
\centering\small
\begin{tabular}{rrrrr}
\toprule
Evaluation dates & Pilot dates & Repeat & Memory ratio & Time ratio\\
\midrule
12,000 & 48,000 & 1 & 0.921 & 0.992\\
12,000 & 48,000 & 2 & 0.937 & 0.989\\
120,000 & 480,000 & 1 & 0.544 & 0.906\\
120,000 & 480,000 & 2 & 0.599 & 1.071\\
\bottomrule
\end{tabular}
\caption{Streaming/dense ratios for complete local pilot tuning, refit and inference. Each route includes imports, input hashing, output persistence and read-back. Input acquisition is excluded equally; resident mapped pages count in peak memory. Order reverses between repeats.}
\label{tab:v10:workload}
\end{table}
Independent dense numerical comparisons cover both covariance choices,
chronological split boundaries and chunks of 1, 5, 37 and 4096 rows. The
largest scaled difference is $1.6\times10^{-14}$ in those checks and
$3.3\times10^{-14}$ in the complete workloads. Audit rank/boundary refusals
and invalid histories retain their meaning. At 120,000 evaluation dates and
480,000 pilot dates, memory ratios are .544 and .599. Elapsed ratios .906 and
1.071 show no consistent speed advantage. The smaller case has only modest
memory differences, because imports and other fixed costs matter. These
four paired comparisons (eight fresh worker runs) establish a local selected memory benefit;
two repeats are too few to characterize a broad performance distribution.
No partial-fit recovery or bank-platform service level follows.

The financial precision account remains material. At the baseline the
fourfold pilot's median interval width is about USD\,498,000; at 4800 dates
it is about USD\,252,000. These values belong to the stipulated 48-branch
synthetic population. They are neither bank estimates nor a stated acceptable
decision width. Extra audit information cannot be assigned a causal cost benefit
from these unpaired design cells; the comparative pilot arms are paired within
each cell. A bank precision target and prices for independent histories and
gold audits are still needed before claiming economic advantage.

